% !TEX root = ../main.tex
\section{Kết luận}
Dự án đã xây dựng một hệ thống ALPR hoàn chỉnh cho biển số Việt Nam, xử lý được cả biển 1 dòng và 2 dòng, đi từ phát hiện, nhận dạng, hậu xử lý theo định dạng biển, theo dõi qua video đến một ứng dụng bãi giữ xe mô phỏng 3D có lấy thẻ, tính phí, so khớp biển ở cổng ra và bảo vệ duyệt các trường hợp bất thường. Backend được tổ chức theo kiến trúc phân lớp với 84 bài kiểm thử tự động.

Bộ phát hiện YOLO11n đạt cả hai chỉ tiêu (mAP@0.5 = 0,995, mAP@0.5:0.95 = 0,713 trên tập test), và hệ thống chạy khoảng 13 FPS trên CPU laptop. Ngược lại, CRNN huấn luyện trên dữ liệu tổng hợp đạt 95,7\,\% trên biển tổng hợp nhưng 0\,\% trên biển thật của bộ dữ liệu. Phân tích lỗi chỉ ra ba nguyên nhân cộng dồn: khoảng cách miền của dữ liệu tổng hợp, phân loại số dòng bằng tỉ lệ khung thất bại trên ảnh bị co giãn, và nhãn chữ chưa được duyệt thật. Các bước tiếp theo theo đúng thứ tự là sửa cách phân loại số dòng, duyệt nhãn, fine-tune CRNN trên biển thật và hoàn thiện đánh giá end-to-end.
